% ═══════════════════════════════════════════════════════════════
% LERNBLATT - DS-01: Mi nueva casa
% Klasse 9, Unidad 2
% ═══════════════════════════════════════════════════════════════

\documentclass[11pt, a4paper]{article}

\usepackage[utf8]{inputenc}
\usepackage[T1]{fontenc}
\usepackage[ngerman]{babel}

% --- SCHRIFTEN ---
\usepackage{helvet}
\renewcommand{\familydefault}{\sfdefault}
\usepackage{palatino}
\newcommand{\seriftext}[1]{{\fontfamily{ppl}\selectfont #1}}

\usepackage{microtype}
\usepackage[margin=1.8cm, top=2cm, bottom=1.5cm]{geometry}
\usepackage{enumitem}
\usepackage{tabularx}
\usepackage{booktabs}
\usepackage{array}
\usepackage{multicol}
\usepackage{tikz}

\usepackage{fancyhdr}
\usepackage{lastpage}

% ═══════════════════════════════════════════════════════════════
% VARIABLEN
% ═══════════════════════════════════════════════════════════════

\newcommand{\lektion}{01}
\newcommand{\thema}{Mi nueva casa}
\newcommand{\untertitel}{Das Haus und der Umzug}
\newcommand{\klasse}{9}
\newcommand{\lernziele}{Hausvokabular · Bewegungsverben · Objektpronomen}

% ═══════════════════════════════════════════════════════════════
% FARBEN
% ═══════════════════════════════════════════════════════════════

\usepackage{xcolor}

\definecolor{colorrojo}{HTML}{C62828}
\definecolor{colorazul}{HTML}{1565C0}
\definecolor{colorverde}{HTML}{2E7D32}
\definecolor{colorgris}{HTML}{757575}
\definecolor{colornegro}{HTML}{222222}

\definecolor{hellgrau}{HTML}{F5F5F5}
\definecolor{rahmen}{HTML}{CCCCCC}
\definecolor{akzent}{HTML}{666666}

% ═══════════════════════════════════════════════════════════════
% BOXEN
% ═══════════════════════════════════════════════════════════════

\usepackage{tcolorbox}
\tcbuselibrary{skins}

\newtcolorbox{lernzielbox}{
    colback=hellgrau, colframe=hellgrau, boxrule=0pt, arc=0pt,
    left=10pt, right=10pt, top=8pt, bottom=8pt,
    before skip=6pt, after skip=8pt
}

\newtcolorbox{grammatikbox}[1][]{
    colback=white, colframe=white, boxrule=0pt, arc=0pt,
    left=10pt, right=6pt, top=6pt, bottom=6pt,
    borderline west={3pt}{0pt}{akzent},
    before skip=4pt, after skip=4pt
}

\newtcolorbox{merkbox}[1][]{
    colback=hellgrau, colframe=hellgrau, boxrule=0pt, arc=0pt,
    left=10pt, right=10pt, top=8pt, bottom=8pt,
    borderline north={2pt}{0pt}{colornegro},
    before skip=6pt, after skip=6pt
}

\newtcolorbox{vokabelbox}[1][]{
    colback=white, colframe=rahmen, boxrule=0.5pt, arc=0pt,
    left=8pt, right=8pt, top=6pt, bottom=6pt,
    before skip=4pt, after skip=4pt
}

\newtcolorbox{tippbox}{
    colback=white, colframe=white, boxrule=0pt, arc=0pt,
    left=8pt, right=4pt, top=3pt, bottom=3pt,
    borderline west={2pt}{0pt}{akzent}
}

% ═══════════════════════════════════════════════════════════════
% BEFEHLE
% ═══════════════════════════════════════════════════════════════

\newcommand{\es}[1]{\seriftext{\textbf{#1}}}
\newcommand{\de}[1]{\textit{\textcolor{akzent}{#1}}}
\newcommand{\gram}[1]{\textsc{#1}}
\newcommand{\abmark}[1]{{\footnotesize\textcolor{akzent}{$\rightarrow$ AB #1}}}
\newcommand{\abschnitt}[2]{\noindent\textbf{\large #1. #2}}
\newcommand{\stamm}[1]{\textcolor{colorrojo}{#1}}

\setlength{\parindent}{0pt}
\setlength{\parskip}{5pt}

% ═══════════════════════════════════════════════════════════════
% KOPF- UND FUSSZEILE
% ═══════════════════════════════════════════════════════════════

\pagestyle{fancy}
\fancyhf{}
\renewcommand{\headrulewidth}{0.5pt}
\renewcommand{\footrulewidth}{0pt}
\lhead{\footnotesize\textsc{Spanisch} · Klasse \klasse}
\rhead{\footnotesize\thema}
\cfoot{\footnotesize\thepage\,/\,\pageref{LastPage}}

% ═══════════════════════════════════════════════════════════════
% DOKUMENT
% ═══════════════════════════════════════════════════════════════

\begin{document}

% --- TITEL ---
\begin{center}
    {\LARGE\bfseries \thema}\\[3pt]
    {\small \untertitel}
\end{center}

\vspace{4pt}

Stell dir vor, du ziehst in eine neue Wohnung in Spanien. Die Umzugshelfer fragen: \es{¿Dónde ponemos el sofá?} -- \de{Wo stellen wir das Sofa hin?} Heute lernst du, wie du dein neues Zuhause beschreiben und einen Umzug organisieren kannst.

\vspace{2pt}

\begin{lernzielbox}
\textbf{Das lernst du:}\quad \lernziele
\end{lernzielbox}

% ═══════════════════════════════════════════════════════════════
% ZWEI SPALTEN
% ═══════════════════════════════════════════════════════════════

\begin{multicols}{2}

% ---------------------------------------------------------------
% ABSCHNITT 1: HAUSRÄUME
% ---------------------------------------------------------------
\abschnitt{1}{Las partes de la casa} \de{-- Die Hausräume}

\vspace{4pt}

Eine Wohnung (\es{el piso}) oder ein Haus (\es{la casa}) hat verschiedene Räume:

\vspace{3pt}

{\small
\begin{tabular}{@{}ll@{}}
\es{el pasillo} & \de{der Flur} \\[0.3em]
\es{el dormitorio} & \de{das Schlafzimmer} \\[0.3em]
\es{el salón} & \de{das Wohnzimmer} \\[0.3em]
\es{la cocina} & \de{die Küche} \\[0.3em]
\es{el balcón} & \de{der Balkon} \\[0.3em]
\es{el baño} & \de{das Badezimmer} \\[0.3em]
\es{el aseo} & \de{die Toilette / Gäste-WC} \\
\end{tabular}}

\vspace{4pt}

\begin{tippbox}
\footnotesize\textit{Tipp:} \es{baño} = Bad mit Dusche/Wanne, \es{aseo} = nur WC.
\end{tippbox}

\abmark{Kasten V1 (oben)}

\vspace{8pt}

% ---------------------------------------------------------------
% ABSCHNITT 2: MÖBEL
% ---------------------------------------------------------------
\abschnitt{2}{Los muebles y objetos} \de{-- Möbel und Gegenstände}

\vspace{4pt}

In den Zimmern stehen verschiedene Möbel (\es{los muebles}):

\vspace{3pt}

{\footnotesize
\begin{tabular}{@{}ll@{}}
\es{la lámpara} & \de{die Lampe} \\
\es{la cama} & \de{das Bett} \\
\es{el armario} & \de{der Schrank} \\
\es{la estantería} & \de{das Regal} \\
\es{la planta} & \de{die Pflanze} \\
\es{el sillón} & \de{der Sessel} \\
\es{el sofá} & \de{das Sofa} \\
\es{el cuadro} & \de{das Bild} \\
\es{el televisor} & \de{der Fernseher} \\
\end{tabular}}

\vspace{3pt}

{\footnotesize
\begin{tabular}{@{}ll@{}}
\es{la nevera} & \de{der Kühlschrank} \\
\es{el cubo de basura} & \de{der Mülleimer} \\
\es{el reloj} & \de{die Uhr} \\
\es{el horno} & \de{der Ofen} \\
\es{la ducha} & \de{die Dusche} \\
\es{el espejo} & \de{der Spiegel} \\
\es{la guitarra} & \de{die Gitarre} \\
\es{el póster} & \de{das Poster} \\
\end{tabular}}

\abmark{Kasten V1 (unten) + Aufgabe 1}

\columnbreak

% ---------------------------------------------------------------
% ABSCHNITT 3: BEWEGUNGSVERBEN
% ---------------------------------------------------------------
\abschnitt{3}{Los verbos de movimiento} \de{-- Bewegungsverben}

\vspace{4pt}

Bei einem Umzug (\es{la mudanza}) brauchst du diese Verben:

\begin{grammatikbox}
{\small
\begin{tabular}{@{}cll@{}}
$\downarrow$ & \es{bajar de} & \de{herunterbringen von} \\[0.4em]
$\uparrow$ & \es{subir a} & \de{hochbringen nach} \\[0.4em]
$\oplus$ & \es{poner en} & \de{stellen/legen in} \\[0.4em]
$\ominus$ & \es{sacar de} & \de{herausnehmen aus} \\[0.4em]
$\leftarrow$ & \es{traer de} & \de{bringen (her) von} \\[0.4em]
$\rightarrow$ & \es{llevar a} & \de{bringen (hin) nach} \\
\end{tabular}}
\end{grammatikbox}

\vspace{3pt}

\begin{merkbox}
\textbf{Merke: Präpositionen}

\vspace{3pt}
{\small
\textbf{de / del} = von, aus (Richtung weg)\\
\textbf{a / al} = nach, zu (Richtung hin)\\
\textbf{en} = in (Position)}
\end{merkbox}

\begin{tippbox}
\footnotesize\textit{Beispiel:} \es{Bajo la caja \textbf{del} camión.} \de{Ich bringe die Kiste aus dem LKW.}
\end{tippbox}

\abmark{Kasten G1 + Aufgabe 2}

\vspace{8pt}

% ---------------------------------------------------------------
% ABSCHNITT 4: OBJEKTPRONOMEN
% ---------------------------------------------------------------
\abschnitt{4}{Los pronombres de OD} \de{-- Objektpronomen}

\vspace{4pt}

Statt das Nomen zu wiederholen, nutzt du Pronomen:

\begin{grammatikbox}
{\small
\begin{tabular}{@{}lll@{}}
& \textbf{Singular} & \textbf{Plural} \\[0.3em]
maskulin & \es{lo} \de{ihn} & \es{los} \de{sie} \\[0.3em]
feminin & \es{la} \de{sie} & \es{las} \de{sie} \\
\end{tabular}}
\end{grammatikbox}

\vspace{3pt}

{\small\textbf{Stellung:} Das Pronomen steht \textbf{vor} dem konjugierten Verb.}

\vspace{3pt}

{\footnotesize
\es{¿La mesa? La ponemos en el salón.}\\
\de{Der Tisch? Wir stellen ihn ins Wohnzimmer.}}

\vspace{3pt}

{\footnotesize
\es{¿Los muebles? Los subimos al piso.}\\
\de{Die Möbel? Wir bringen sie in die Wohnung hoch.}}

\abmark{Kasten G2 + Aufgabe 3}

\end{multicols}

% ═══════════════════════════════════════════════════════════════
% REDEMITTEL
% ═══════════════════════════════════════════════════════════════

\vspace{4pt}

\abschnitt{5}{Comunicación} \de{-- Der Umzugsdialog}

\vspace{4pt}

\begin{vokabelbox}
\begin{multicols}{2}

{\small\textbf{Fragen:}}

\vspace{2pt}

{\footnotesize
\es{¿Dónde ponemos el/la...?}\\
\de{Wo stellen wir den/die... hin?}

\vspace{3pt}

\es{¿Dónde está el/la...?}\\
\de{Wo ist der/die...?}

\vspace{3pt}

\es{¿Qué hay en...?}\\
\de{Was gibt es in...?}}

\columnbreak

{\small\textbf{Antworten:}}

\vspace{2pt}

{\footnotesize
\es{Lo/La ponemos en el/la...}\\
\de{Wir stellen ihn/sie in den/die...}

\vspace{3pt}

\es{Está al lado de / detrás de / en...}\\
\de{Er/Sie ist neben / hinter / in...}

\vspace{3pt}

\es{Hay un/una...}\\
\de{Es gibt einen/eine...}}

\end{multicols}
\end{vokabelbox}

\abmark{Kasten R1 + Aufgabe 4}

\vspace{8pt}

% ═══════════════════════════════════════════════════════════════
% ZUSAMMENFASSUNG
% ═══════════════════════════════════════════════════════════════

\begin{merkbox}
\textbf{Zusammenfassung: Was du heute gelernt hast}

\vspace{4pt}

\begin{multicols}{3}
{\footnotesize

\textbf{Vokabeln:}\\
7 Hausräume\\
17 Möbel/Gegenstände

\columnbreak

\textbf{Grammatik:}\\
6 Bewegungsverben + Präp.\\
4 Objektpronomen (lo/la/los/las)

\columnbreak

\textbf{Kommunikation:}\\
Umzug beschreiben\\
Fragen + Antworten

}
\end{multicols}
\end{merkbox}

\end{document}
